% 00-portada-resumen.tex — Portada, resúmenes ES/EN, metadatos de la entrega.
% Fuente: docs/informe-final (borrador previo), verificado contra
% CONTRIBUTORS.md, README.md, CITATION.cff y docs/observaciones/.

\begin{titlepage}
\centering

\textsc{\Large Universidad Técnica Estatal de Quevedo}\\[0.3cm]
\textsc{\large Facultad de Ciencias de la Computación y Diseño Digital}\\
\textsc{\large Carrera de Ingeniería de Software}\\[1.5cm]

\rule{\linewidth}{0.5mm}\\[0.4cm]
{\LARGE \bfseries ARTISYNC: Plataforma web de intermediación entre
creadores de contenido digital y clientes con estrategia híbrida de acceso
a datos, verificación asistida por IA y evaluación empírica
multidimensional}\\[0.4cm]
\rule{\linewidth}{0.5mm}\\[1cm]

{\large \gls{pfc} --- Aplicaciones Web, Quinto Nivel}\\
{\large Período Académico Presencial 2026-2027}\\[1.5cm]

\begin{minipage}{0.9\textwidth}
\centering
\begin{tabular}{@{}ll@{}}
\toprule
\textbf{Integrante} & \textbf{ORCID} \\
\midrule
Bone Arroyo, Niurca Scarleth & \url{https://orcid.org/0009-0002-2219-2800} \\
Carvajal Loor, Johan Stalin & \url{https://orcid.org/0009-0008-9229-382X} \\
Figueroa Morales, Bryan Javier & \url{https://orcid.org/0009-0009-6357-4996} \\
Rios Cuyabazo, Jhon Kevin & \url{https://orcid.org/0009-0003-7446-9450} \\
\bottomrule
\end{tabular}
\end{minipage}\\[1cm]

\textbf{Docente-director:} Dr. Gleiston Cicerón Guerrero Ulloa, Ph.D.\\
\texttt{gguerrero@uteq.edu.ec}\\[1cm]

\begin{minipage}{0.85\textwidth}
\centering
\textbf{Fecha de cierre:} 11 de septiembre de 2026 (semana 19)\\
\textbf{Etiqueta del artefacto:} \texttt{v1.0.0}\\
\textbf{DOI del software (Zenodo):} \texttt{10.5281/zenodo.21978572}
(archivo de \texttt{v1.0.0}; la versión anterior, \texttt{v0.9.0-rc},
quedó archivada aparte en \texttt{10.5281/zenodo.21730559})\\
\textbf{DOI del dataset de mediciones (Zenodo):}
\texttt{10.5281/zenodo.22236251} (depósito separado del software,
licencia CC~BY~4.0, conforme al Bloque D.3 de la guía)
\end{minipage}

\vfill
{\small Repositorio: \url{https://github.com/JohanCarvajal04/Proyecto-WEB-ARTISYNC}}
\end{titlepage}

\chapter*{Resumen}
\addcontentsline{toc}{chapter}{Resumen}

\textbf{Contexto y problema.} La economía creativa digital fragmenta la contratación, comunicación y pagos entre plataformas desconectadas, aumentando el riesgo de fraude. \textbf{Objetivo.} Diseñar, implementar y evaluar Artisync, plataforma construida sobre Spring Boot y Angular que integra el ciclo completo de intermediación con pagos \emph{escrow}, usando una estrategia híbrida de acceso a datos (\gls{orm} y PL/pgSQL) y verificación por IA. \textbf{Métodos.} Se aplicó la metodología \gls{dsr} de Peffers~\cite{PEFFERS2007}, ingeniería de requisitos ISO/IEC/IEEE~29148:2018~\cite{ISO29148}, y evaluación empírica de rendimiento (k6), seguridad (controles OWASP Top 10 y análisis dinámico con ZAP), usabilidad (\gls{sus}, $N=16$)~\cite{BROOKE1996,BANGOR2009}, cobertura (JaCoCo) y calidad web (Lighthouse), sobre un entorno reproducible con Docker Compose, con comparaciones de rendimiento validadas mediante pruebas inferenciales no paramétricas y tamaño de efecto ordinal. \textbf{Resultados principales.} Artisync logró p95 de 50.17~ms bajo carga, 0\,\% de errores HTTP, cobertura JaCoCo de 82.93\,\% líneas y 71.50\,\% ramas ---las tres capas sobre el umbral de 70\,\%---, Lighthouse Performance 100/100 (desktop) y 80 (mobile), y 0 hallazgos altos en ZAP. La usabilidad, en cambio, alcanzó un \gls{sus} de 61.25/100, sin superar el umbral de aceptabilidad (68). Se implementaron 26 procedimientos almacenados sin SQL dinámico, y los 37 requisitos declarados conservan trazabilidad automática hacia código y pruebas, validada en \gls{ci}. \textbf{Conclusiones.} Artisync demuestra la viabilidad de la estrategia híbrida y de la validación asistida por IA en 17 semanas, con evidencia reproducible; la usabilidad es la única dimensión que no alcanza su umbral y concentra el trabajo futuro.

\textbf{Palabras clave:} ingeniería de requisitos; procedimientos
almacenados; verificación de identidad; arquitectura de microservicios;
usabilidad de software; seguridad web

\chapter*{Abstract}
\addcontentsline{toc}{chapter}{Abstract}

\textbf{Context and problem.} The digital creative economy fragments contracting, communication, and payments across disconnected platforms, increasing fraud risks. \textbf{Objective.} To design, implement, and evaluate Artisync, a platform built on Spring Boot and Angular integrating the intermediation cycle with escrow payments, a hybrid data-access strategy (ORM and PL/pgSQL), and AI-assisted verification. \textbf{Methods.} Peffers' \gls{dsr} methodology~\cite{PEFFERS2007} was applied, alongside ISO/IEC/IEEE~29148:2018 requirements engineering~\cite{ISO29148}, and empirical evaluations of performance (k6), security (OWASP Top 10 controls and dynamic scanning with ZAP), usability (\gls{sus}, $N=16$)~\cite{BROOKE1996,BANGOR2009}, code coverage (JaCoCo), and web quality (Lighthouse), on an environment reproducible with Docker Compose, with performance comparisons validated through non-parametric inferential tests and ordinal effect sizes. \textbf{Main results.} Artisync achieved a p95 of 50.17~ms under load, 0\% HTTP errors, JaCoCo coverage of 82.93\% lines and 71.50\% branches ---all three layers above the 70\% threshold---, Lighthouse Performance 100/100 (desktop) and 80 (mobile), and 0 high security findings in ZAP. Usability, by contrast, reached a \gls{sus} of 61.25/100, failing to pass the acceptability threshold (68). 26 stored procedures free of dynamic SQL were implemented, and the 37 declared requirements maintain automated traceability to code and tests, validated in \gls{ci}. \textbf{Conclusions.} Artisync proves the viability of hybrid data access and AI-assisted validation in a 17-week timeframe, with reproducible evidence; usability is the only dimension falling short of its threshold and concentrates the future work.

\textbf{Keywords:} requirements engineering; stored procedures; identity
verification; microservices architecture; software usability; web
security
